\documentclass[11pt]{article}
\usepackage[margin=1in]{geometry}
\usepackage{amsmath,amssymb,amsthm}
\usepackage{booktabs}
\usepackage{hyperref}

\newtheorem{theorem}{Theorem}
\newtheorem{lemma}[theorem]{Lemma}
\newtheorem{proposition}[theorem]{Proposition}
\newtheorem{conjecture}[theorem]{Conjecture}
\newtheorem{definition}[theorem]{Definition}
\newtheorem{corollary}[theorem]{Corollary}
\newtheorem{remark}[theorem]{Remark}

\title{Prime Geometry I:\\
The Prime Triangle and its Geometric Invariants}
\author{Allen Proxmire}
\date{April 2026}

\begin{document}
\maketitle

\begin{abstract}
To each pair of consecutive primes $(p_n, p_{n+1})$ we associate a
right triangle with legs $p_n$ and $p_{n+1}$. The angle
$\alpha_n = \arctan(p_n/p_{n+1})$ and its associated monotone proxy
$\rho_n = p_n / p_{n+1}$ are the central geometric quantities; the
hypotenuse $C_n = \sqrt{p_n^{\,2} + p_{n+1}^{\,2}}$ and, for triples,
the difference $C_2 - C_1$ generate a family of elementary
invariants. We derive the Prime Square-Difference identity
$C_2^{\,2} - C_1^{\,2} = p_{n+2}^{\,2} - p_n^{\,2}$, define the
integer-valued PSD factor, and record first-order approximations
for the discrete energy $E_n$, curvature $K_n$, angle drift
$\Delta \alpha_n$, and normalized shape curvature $\chi_n$. We then
introduce the angle-record sequence and state the $2P$-beats lemma
that governs it, foreshadowing the dyadic-density interpretation of
PG~II and the Generalized Bertrand Principle of PG~III. Finally we
position this framework within the Ramanujan-prime family, from
which the twin- and constellation-Ramanujan sequences arise as
immediate analogs.
\end{abstract}

\section{Introduction}

Prime gaps exhibit rich structure that is invisible at the level of
the primes alone. This paper introduces a minimal geometric
construction which encodes consecutive primes as right triangles and
records prime-gap dynamics in the associated angles. The construction
is elementary, and most of its identities are algebraically routine.
Its value lies not in depth but in \emph{angle}: by shifting the
viewpoint from gaps to triangles, one discovers quantities whose
record structure encodes non-trivial density information about the
primes.

The Prime Triangle framework can be thought of as a \emph{geometric
microscope} on prime gaps. Consecutive-prime pairs produce
triangles whose angles cluster near $45^\circ$, with twin primes
approaching the limit fastest. Asking which pairs set new records in
the angle sequence turns out to be equivalent to asking about the
density of twin primes in dyadic intervals~--- an observation
developed fully in PG~II and generalized in PG~III. The present note
establishes the geometric objects and identities on which those
companion papers rest.

Three themes run through the paper:
\begin{enumerate}
\item[(1)] \emph{Geometric encoding.} The triangle
$(p_n, p_{n+1}, C_n)$ records a prime pair; the angle
$\alpha_n = \arctan(p_n/p_{n+1})$ records the asymmetry between the
pair's members.
\item[(2)] \emph{Derived quantities.} Elementary algebra yields an
exact identity $C_2^{\,2} - C_1^{\,2} = p_{n+2}^{\,2} - p_n^{\,2}$
and an always-integer quantity $\mathrm{PSD}_n$. Gap-scaled
approximations produce a discrete energy $E_n$, curvature $K_n$, and
normalized shape curvature $\chi_n$.
\item[(3)] \emph{Dyadic interpretation.} The sequence of
angle-records among consecutive prime pairs satisfies an
elementary comparison rule that translates directly into a dyadic
density statement on twin primes and, more generally, on admissible
prime constellations. This is the hinge between the geometric
object (the Prime Triangle) and the arithmetic statements of
PG~II--III.
\end{enumerate}

Throughout, $\mathbb{P}$ denotes the primes, $(p_n)_{n\ge 1}$ the
ordered list of primes (so $p_1 = 2$, $p_2 = 3$, $\ldots$), and
$\pi(x) = \#\{p \in \mathbb{P} : p \le x\}$ the prime-counting
function. For pair-constellation $\mathcal{C} = \{0, g\}$ with
$g \ge 2$ even, the counting function
$\pi_{\mathcal{C}}(x) = \#\{p \le x : p \in \mathbb{P},\;p+g \in \mathbb{P}\}$
specializes to $\pi_2(x)$ when $g = 2$.

\section{The Prime Triangle}

\subsection{Construction}

\begin{definition}[Prime Triangle]
For consecutive primes $p_n < p_{n+1}$, the \emph{Prime Triangle}
at index $n$ is the right triangle with legs $p_n$ and $p_{n+1}$
and hypotenuse
\[
C_n \;:=\; \sqrt{p_n^{\,2} + p_{n+1}^{\,2}}.
\]
We record the pair's geometric parameters as
\[
\alpha_n \;:=\; \arctan\!\left(\frac{p_n}{p_{n+1}}\right),
\qquad
\rho_n \;:=\; \frac{p_n}{p_{n+1}}.
\]
The quantity $\alpha_n$ is the angle at the base of the leg of
length $p_{n+1}$, and $\rho_n = \tan \alpha_n$.
\end{definition}

\begin{proposition}[Basic properties]\label{prop:basic}
For every $n\ge 1$:
\begin{enumerate}
\item[(i)] $0 < \alpha_n < 45^\circ$, equivalently $0 < \rho_n < 1$.
\item[(ii)] $\rho$ and $\alpha$ are strictly monotone: for pairs
$(p,q)$ and $(p',q')$ with $p<q$ and $p'<q'$,
$\rho(p,q) > \rho(p',q') \iff \alpha(p,q) > \alpha(p',q')$.
\item[(iii)] For a twin pair $(p, p+2)$,
$\rho(p,p+2) = p/(p+2)$ and $\alpha(p,p+2) \to 45^\circ$ as
$p \to \infty$.
\item[(iv)] For any fixed gap $g \ge 2$, $\alpha(p, p+g) \to 45^\circ$
as $p \to \infty$; twin pairs achieve this limit fastest, followed
by cousin pairs, sexy pairs, and so on.
\end{enumerate}
\end{proposition}

\begin{proof}
(i) follows from $p_n < p_{n+1}$, giving $\rho_n < 1$ and hence
$\alpha_n < 45^\circ$. (ii) is immediate since $\arctan$ is strictly
increasing on $\mathbb{R}_{>0}$. (iii) is direct computation: for twins,
$\rho(p,p+2) = 1 - 2/(p+2) \to 1$. (iv) For gap $g$,
$\rho(p,p+g) = 1 - g/(p+g)$; for fixed $g$ and large $p$,
$1 - \rho(p,p+g) \sim g/p$, so smaller $g$ gives faster convergence.
\end{proof}

\subsection{Numerical illustration}

The first few Prime Triangles:

\begin{center}
\begin{tabular}{rrrrr}
\toprule
$n$ & $p_n$ & $p_{n+1}$ & gap & $\alpha_n$ \\
\midrule
1 & 2 & 3 & 1 & $33.690^\circ$ \\
2 & 3 & 5 & 2 & $30.964^\circ$ \\
3 & 5 & 7 & 2 & $35.538^\circ$ \\
4 & 7 & 11 & 4 & $32.471^\circ$ \\
5 & 11 & 13 & 2 & $40.236^\circ$ \\
6 & 13 & 17 & 4 & $37.405^\circ$ \\
7 & 17 & 19 & 2 & $41.820^\circ$ \\
8 & 19 & 23 & 4 & $39.558^\circ$ \\
9 & 23 & 29 & 6 & $38.428^\circ$ \\
10 & 29 & 31 & 2 & $43.091^\circ$ \\
\bottomrule
\end{tabular}
\end{center}

Observe the behavior: twin pairs consistently produce the largest
angles at each scale, with the angle monotonically approaching
$45^\circ$ along the twin subsequence. Non-twin pairs interleave
between twin records and always fall short of the preceding twin
angle. This empirical pattern is made precise by the $2P$-beats lemma
(Section~\ref{sec:angle-records}) and the angle-record theorem of
PG~II.

\section{Derived Quantities}

The Prime Triangle construction extends naturally to triples of
consecutive primes, producing the identities and derived invariants
collected in this section. Throughout, fix three consecutive primes
$p_n < p_{n+1} < p_{n+2}$ and write
\[
C_1 \;=\; \sqrt{p_n^{\,2} + p_{n+1}^{\,2}},
\qquad
C_2 \;=\; \sqrt{p_{n+1}^{\,2} + p_{n+2}^{\,2}}.
\]

\subsection{The Prime Square-Difference Identity}

\begin{theorem}[PSD identity]\label{thm:psd}
For every triple of consecutive primes,
\[
C_2^{\,2} - C_1^{\,2} \;=\; p_{n+2}^{\,2} - p_n^{\,2}.
\]
Equivalently, writing $G_n := p_{n+2} - p_n$ for the skip-one gap,
\[
C_2^{\,2} - C_1^{\,2} \;=\; G_n\,(2 p_n + G_n).
\]
\end{theorem}

\begin{proof}
Direct expansion:
\[
C_2^{\,2} - C_1^{\,2}
= \big[p_{n+2}^{\,2} + p_{n+1}^{\,2}\big]
- \big[p_n^{\,2} + p_{n+1}^{\,2}\big]
= p_{n+2}^{\,2} - p_n^{\,2}
= (p_{n+2} - p_n)(p_{n+2} + p_n)
= G_n(2p_n + G_n). \qedhere
\]
\end{proof}

\begin{definition}[PSD factor]
The \emph{Prime Square-Difference factor} at index $n$ is
\[
\mathrm{PSD}_n \;:=\; \frac{p_{n+2}^{\,2} - p_n^{\,2}}{12}
\;=\; \frac{G_n\,(2 p_n + G_n)}{12}.
\]
\end{definition}

\begin{theorem}[Integrality and last-digit structure]\label{thm:psd-int}
For every $n$ with $p_n \ge 5$:
\begin{enumerate}
\item[(i)] $\mathrm{PSD}_n \in \mathbb{Z}$.
\item[(ii)] The last decimal digit of $\mathrm{PSD}_n$ is one of
$\{0, 4, 6\}$.
\end{enumerate}
\end{theorem}

\begin{proof}
Every prime $p \ge 5$ satisfies $p \equiv \pm 1 \pmod 6$. Writing
$p_n = 6a \pm 1$, $p_{n+2} = 6b \pm 1$, a four-case expansion gives
\[
(6b \pm 1)^{2} - (6a \pm 1)^{2} = 12(b - a)\big[3(a+b) \pm 1\big],
\]
so $12 \mid p_{n+2}^{\,2} - p_n^{\,2}$ and $\mathrm{PSD}_n$ is an
integer. For (ii), reducing modulo $120 = \mathrm{lcm}(12, 10)$ gives
$p_{n+2}^{\,2} - p_n^{\,2} \in \{0, 48, 72\} \pmod{120}$, so
$\mathrm{PSD}_n \in \{0, 4, 6\} \pmod{10}$.
\end{proof}

\begin{remark}[Twin-structure specialization]
\label{rem:psd-twin}
When $G_n = 6$, the three consecutive primes take one of the forms
$(p_n, p_n+2, p_n+6)$ or $(p_n, p_n+4, p_n+6)$; in either case
$p_{n+1}$ is a member of a twin prime pair. A direct computation
from the identity gives
\[
\mathrm{PSD}_n \;=\; \frac{6(2p_n + 6)}{12}
\;=\; \frac{p_n + p_{n+2}}{2} \;=\; p_{n+1} \pm 1,
\]
where the sign encodes which adjacent pair forms the twin. This is
the simplest case in which $\mathrm{PSD}_n$ pinpoints twin-prime
structure within a triple.
\end{remark}

\subsection{Energy, curvature, and angle drift}

\begin{definition}[Energy and skip-one gap]
For three consecutive primes, define the \emph{energy transition}
\[
E_n \;:=\; C_2 - C_1,
\]
and recall $G_n = p_{n+2} - p_n = g_n + g_{n+1}$ with $g_n = p_{n+1} - p_n$.
\end{definition}

\begin{proposition}[First-order expansion of $E_n$]\label{prop:energy}
With $g_n, g_{n+1} \ll p_n$,
\[
E_n \;=\; \frac{\sqrt{2}}{2}\,G_n \;+\; O\!\left(\frac{G_n^{\,2}}{p_n}\right),
\]
and more precisely
\[
\frac{E_n}{G_n} \;\approx\; \frac{\sqrt{2}}{2}
- \frac{g_n - g_{n+1}}{2\sqrt{2}\,(2p_n + G_n)}.
\]
\end{proposition}

\begin{proof}[Proof sketch]
Taylor-expand $\sqrt{(p_{n+1}+g_{n+1})^2 + p_{n+1}^{\,2}}
- \sqrt{p_n^{\,2} + p_{n+1}^{\,2}}$ in $g_{n+1}/p_{n+1}$ and
$g_n/p_n$ and collect first-order terms.
\end{proof}

\begin{definition}[Discrete curvature and shape curvature]
\begin{align*}
K_n &:= E_{n+1} - E_n \qquad \text{(discrete curvature)},\\
\chi_n &:= K_n / E_n \qquad \text{(normalized shape curvature)}.
\end{align*}
\end{definition}

\begin{proposition}\label{prop:chi}
Under the same first-order approximation,
\[
K_n \;\approx\; \frac{\sqrt{2}}{2}\,(g_{n+2} - g_n),
\qquad
\chi_n \;\approx\; \frac{g_{n+2} - g_n}{g_n + g_{n+1}}.
\]
\end{proposition}

The quantity $\chi_n$ is a dimensionless, scale-free measure of how
the prime-gap sequence bends at index $n$. It will be the central
object of PG~II and PG~III's empirical sections when examining
sign-coherence and long-range structure; for the present paper it is
enough to record its definition.

\subsection{Angle drift}

\begin{definition}[Angle drift]
\[
\Delta \alpha_n \;:=\; \alpha_{n+1} - \alpha_n.
\]
\end{definition}

\begin{proposition}[Derivative hierarchy]\label{prop:hierarchy}
In the first-order approximation, $\Delta \alpha_n$ behaves as a
discrete first derivative of the sequence $\{\alpha_n\}$, and
$\chi_n$ as a normalized discrete second derivative. Schematically,
\[
\chi_n \;\longrightarrow\; \Delta \alpha_n \;\longrightarrow\; \alpha_n,
\]
with each arrow indicating discrete integration.
\end{proposition}

The hierarchy explains, in an elementary way, why the angle sequence
$\alpha_n$ exhibits long coherent phases and suppressed extremes:
curvature $\chi_n$ is small on average, so $\Delta\alpha_n$ is
well-controlled, so $\alpha_n$ drifts smoothly toward its $45^\circ$
limit.

\section{Angle Dynamics and Records}\label{sec:angle-records}

\subsection{The angle-record sequence}

\begin{definition}[Angle-record]
A consecutive prime pair $(p_n, p_{n+1})$ is an
\emph{angle-record} if
$\rho_n > \rho_m$ for every $m < n$, equivalently
$\alpha_n > \alpha_m$ for every $m < n$.
\end{definition}

By Proposition~\ref{prop:basic}, the angle-record sequence is an
infinite subsequence of the consecutive-prime-pair sequence. Its
first few terms are:

\begin{center}
\begin{tabular}{rrrrr}
\toprule
record \# & $p_n$ & $p_{n+1}$ & $\alpha_n$ & type \\
\midrule
1 & 2 & 3 & $33.690^\circ$ & gap-1 initial pair \\
2 & 5 & 7 & $35.538^\circ$ & twin \\
3 & 11 & 13 & $40.236^\circ$ & twin \\
4 & 17 & 19 & $41.820^\circ$ & twin \\
5 & 29 & 31 & $43.091^\circ$ & twin \\
6 & 41 & 43 & $43.636^\circ$ & twin \\
7 & 59 & 61 & $44.045^\circ$ & twin \\
8 & 71 & 73 & $44.204^\circ$ & twin \\
$\vdots$ & $\vdots$ & $\vdots$ & $\vdots$ & $\vdots$ \\
\bottomrule
\end{tabular}
\end{center}

\subsection{The \texorpdfstring{$2P$}{2P}-beats lemma}

The key elementary observation governing the record structure is
the following comparison rule.

\begin{lemma}[$2P$-beats]\label{lem:2P}
Let $(p, p+g)$ and $(q, q+h)$ be prime pairs with $g, h \ge 2$. Then
\[
\rho(p, p+g) \;>\; \rho(q, q+h)
\quad\Longleftrightarrow\quad
h\, p \;>\; g\, q.
\]
In particular, a gap-$g$ pair at $p$ beats a twin pair $(q, q+2)$ in
angle if and only if $p > (g/2)\, q$; the binding case (smallest $g$
larger than $2$) is $g = 4$, requiring $p > 2q$.
\end{lemma}

\begin{proof}
$\rho(p, p+g) > \rho(q, q+h)$
$\iff p(q+h) > q(p+g)$
$\iff hp > gq$.
\end{proof}

\begin{corollary}[Preview of the angle-record theorem]
\label{cor:angle-record-preview}
Let $(T_k)_{k\ge 1}$ be the twin-prime subsequence
(smaller members). Suppose that for every $k$ with $T_k \ge 11$,
$T_{k+1} < 2 T_k$. Then every angle-record with leading prime
$p_n \ge 3$ is a twin prime.
\end{corollary}

\begin{proof}[Proof sketch]
If a non-twin pair $(p, p+g)$ with $g \ge 4$ were to set a record
after the twin $T_k$ and before $T_{k+1}$, Lemma~\ref{lem:2P} would
require $p > (g/2) T_k \ge 2 T_k$. But $p \le T_{k+1} < 2 T_k$ by
hypothesis; contradiction.
\end{proof}

The converse direction and the full equivalence
with the dyadic density inequality $\pi_2(2x) - \pi_2(x) \ge 1$ are
developed in PG~II.

\section{Geometric Interpretation of Dyadic Density}

The $2P$-beats lemma reveals a direct equivalence between angle
dynamics and dyadic density. Corollary~\ref{cor:angle-record-preview}
asserts that the factor-$2$ bound on twin-prime ratios
$T_{k+1} < 2 T_k$ is exactly what guarantees twins dominate the
angle-record sequence. Conversely, if the angle-record sequence
omits a twin prime (i.e., some non-twin pair sets a record), then
some ratio $T_{k+1}/T_k \ge 2$ must occur.

\subsection{Bertrand factors, revisited geometrically}

The dyadic factor $2$ is, in classical terms, Bertrand's factor:
every integer $n > 1$ has a prime in $(n, 2n]$. For twin primes, the
analogous dyadic statement
\[
\pi_2(2x) - \pi_2(x) \;\ge\; 1 \qquad (x \ge 11),
\]
which we formalize as the Twin-Prime Bertrand Postulate
$(\mathrm{TPB})$ in PG~II, is exactly the ratio condition
$T_{k+1} < 2 T_k$ for $T_k \ge 11$ via a direct one-line argument.

Geometrically, Lemma~\ref{lem:2P} explains the recurrence of the
factor $2$: for a gap-$g$ constellation $\mathcal{C} = \{0, g\}$ and
its \emph{doubled companion} $\mathcal{C}' = \{0, 2g\}$, a
$\mathcal{C}'$-pair at $P$ beats a $\mathcal{C}$-pair at $Q$ in
angle if and only if $P > 2Q$ --- independent of $g$. The dyadic
factor $2$ is therefore also the universal angle-beats factor
between doubled constellations, providing a geometric rationale for
the uniform Bertrand factor that appears in PG~III's Generalized
Bertrand Principle.

\subsection{A summary of the geometric--arithmetic correspondence}

We can summarize the correspondence in one line:

\begin{center}
\emph{Angle-records dominated by a gap-$g$ constellation~~$\iff$~~dyadic
density $\pi_{\mathcal{C}}(2x) - \pi_{\mathcal{C}}(x) \ge 1$ holds
for all $x$ past a threshold.}
\end{center}

For $g = 2$ (twins), the dominance is empirically total once we exclude
the initial non-twin $(2,3)$; the equivalent dyadic inequality is
$(\mathrm{TPB})$. For larger $g$ the statement is modified by the
role of smaller-gap competitors (twins beat cousins at angle
almost everywhere), but the \emph{within-constellation} dyadic
inequality
\[
\pi_{\mathcal{C}}(2x) - \pi_{\mathcal{C}}(x) \;\ge\; 1
\]
persists with the same Bertrand factor of $2$. The full picture
is developed in PG~III as the Generalized Bertrand Principle.

\section{The Ramanujan-Prime Context}

Bertrand's postulate $\pi(2x) - \pi(x) \ge 1$ for $x \ge 1$,
proved by Chebyshev in 1852, is the $n = 1$ case of a broader family.
Ramanujan~\cite{Ramanujan1919} proved in 1919 that for every integer
$n \ge 1$, the inequality
$\pi(x) - \pi(x/2) \ge n$ holds for all sufficiently large $x$; the
$n$-th \emph{Ramanujan prime} is defined as
\begin{equation}\label{eq:Rn-classical}
R_n \;=\; \min\Big\{\,R\,:\,\pi(x) - \pi(x/2) \ge n\ \text{for all}\ x \ge R\,\Big\},
\end{equation}
so $R_1 = 2$. The first several values are $R_n = 2, 11, 17, 29,
41, 59, 67, 71, \ldots$ (OEIS \texttt{A104272}~\cite{OEIS_A104272}).
Sondow~\cite{Sondow2009} established the asymptotic
$R_n \sim p_{2n}$ and initiated the modern study of this sequence;
later work~\cite{Amersi2014, Paksoy2012} introduced
$c$-Ramanujan primes (for interval ratios other than $1/2$) and
derived Ramanujan-prime iterations.

The Prime Triangle framework makes two direct analogs transparent:
\begin{itemize}
\item \emph{Twin-Ramanujan primes.} Replacing $\pi$ by $\pi_2$ in
\eqref{eq:Rn-classical} gives the sequence
\[
R^{\mathrm{twin}}_n \;=\; \min\Big\{\,R\,:\,\pi_2(x) - \pi_2(x/2) \ge n\ \text{for all}\ x \ge R\,\Big\}.
\]
Here $R^{\mathrm{twin}}_1$ is exactly the Twin-Prime Bertrand
threshold $(\mathrm{TPB})$, equal to $11$. Higher values
$R^{\mathrm{twin}}_n$ quantify the higher-order dyadic density of
twin primes.
\item \emph{Constellation-Ramanujan primes.} For each admissible
$\mathcal{C} = \{0, g\}$ the sequence
\[
R^{\mathcal{C}}_n \;=\; \min\Big\{\,R\,:\,\pi_{\mathcal{C}}(x) - \pi_{\mathcal{C}}(x/2) \ge n\ \text{for all}\ x \ge R\,\Big\}
\]
is the natural extension; the cases $g = 4, 6$ give
$R^{\{0,4\}}_1 = 7$ and $R^{\{0,6\}}_1 = 5$.
\end{itemize}

Each of these sequences is empirically a subsequence of its
defining constellation (every $R^{\mathrm{twin}}_n$ is a twin prime,
every $R^{\mathcal{C}}_n$ is a member of $\mathcal{C}$), because the
relevant counting function jumps upward only at constellation
members. PG~II develops $R^{\mathrm{twin}}_n$ explicitly and PG~III
generalizes to other $\mathcal{C}$; the present paper provides the
geometric setting in which these are natural.

\section{Conclusion}

The Prime Triangle construction is elementary. Its value is in how
it organizes. The angle $\alpha_n$ is monotone-equivalent to the
ratio $\rho_n = p_n / p_{n+1}$, both of which record the geometric
shape of consecutive primes. The triangle's hypotenuse difference
$C_2 - C_1$ yields an exact identity with the skip-one square
difference $p_{n+2}^{\,2} - p_n^{\,2}$, producing the integer-valued
$\mathrm{PSD}_n$ factor. First-order expansions give an energy
$E_n$, curvature $K_n$, normalized shape curvature $\chi_n$, and
angle drift $\Delta\alpha_n$ that form a derivative hierarchy.

These objects are the geometric setting for the main conjectural
content of the PG series:

\begin{itemize}
\item The sequence of angle-records selects, via the $2P$-beats
lemma, a subfamily of consecutive-prime pairs that is equivalent
(up to the trivial $(2, 3)$) to the twin primes. This equivalence
is the Angle-Record Theorem of PG~II.
\item The equivalence extends via the same $2P$-beats threshold to
a dyadic density statement: at least one twin prime in every
$(x, 2x]$ with $x$ past an effective threshold. This is the
Twin-Prime Bertrand Postulate of PG~II, with threshold
$R^{\mathrm{twin}}_1 = 11$.
\item Generalization to admissible pair-constellations yields the
Generalized Bertrand Principle of PG~III, and a corresponding
family of constellation-Ramanujan primes $R^{\mathcal{C}}_n$ that
extends the Ramanujan-prime framework beyond ordinary primes for
the first time.
\end{itemize}

The Prime Triangle is thus a small object with a long reach: from
an elementary algebraic identity through an elementary geometric
comparison lemma to an explicit, unproven, well-posed conjectural
statement on prime-constellation density. The arithmetic
consequences are developed in the companion papers.

\section*{Acknowledgements}

This paper is the first in a three-part series. PG~II introduces
the Twin-Prime Bertrand Postulate and proves the Angle-Record
Theorem. PG~III generalizes to admissible pair-constellations and
establishes the Generalized Bertrand Principle, including the
twin- and constellation-Ramanujan prime sequences.

\begin{thebibliography}{99}

\bibitem{Ramanujan1919}
S.~Ramanujan.
\emph{A proof of Bertrand's postulate}.
J.~Indian Math.\ Soc.\ \textbf{11} (1919), 181--182.

\bibitem{Sondow2009}
J.~Sondow.
\emph{Ramanujan primes and Bertrand's postulate}.
Amer.\ Math.\ Monthly \textbf{116} (2009), 630--635.
arXiv:\href{https://arxiv.org/abs/0907.5232}{0907.5232}.

\bibitem{Amersi2014}
N.~Amersi, O.~Beckwith, S.~J.~Miller, R.~Ronan, J.~Sondow.
\emph{Generalized Ramanujan primes}.
In: \emph{Combinatorial and Additive Number Theory~--- CANT 2011 and 2012},
Springer Proc.\ Math.\ Stat.\ \textbf{101} (2014), 1--13.

\bibitem{Paksoy2012}
M.~B.~Paksoy.
\emph{Derived Ramanujan primes: $R'_n$}.
arXiv:\href{https://arxiv.org/abs/1210.6991}{1210.6991} (2012).

\bibitem{OEIS_A104272}
OEIS Foundation Inc.
\emph{Entry A104272: Ramanujan primes}.
Online at \url{https://oeis.org/A104272}.

\end{thebibliography}

\end{document}
